% GENERATED FILE. Edit canonical lessons, then run npm run generate:books.
% canonical-source-hash: fnv1a64:d1dedae9a35b2326
% canonical-lessons: GE-C324-ruehren, GE-C324-umruehren, GE-C324-verquirlen, GE-C324-aufschlagen, GE-C324-schneebesen

\chapter{Rühren}
\label{ch:ge-ruhren}
\hlchaptermodality{324}

\begin{chapteropening}
\textbf{By the end of this chapter:} \emph{I can describe mixing: to stir, to give a stir, to whisk, to crack an egg and a whisk.}

Complete the last lesson of chapter 324: I can describe mixing: to stir, to give a stir, to whisk, to crack an egg and a whisk.
\end{chapteropening}

\section[rühren]{rühren — to stir; to move (the heart)}

\label{lesson:GE-C324-ruehren}



\begin{quote}
Before the new one: say the German for to grate, then the German for to halve.
\end{quote}

\subsection*{You'll want to know: rühren}
\textbf{rühren} — \textquotedblleft{}to stir; to move (the heart)\textquotedblright{}.

\begin{quote}
Er rührte die Suppe langsam.
\end{quote}

\begin{quote}
\emph{He stirred the soup slowly.}
\end{quote}

Also to touch the heart: \emph{Das Lied rührte sie}.

\begin{cousinweb}
\emph{rühren} is inherited from Old High German \emph{ruoren}, from Proto-Germanic \emph{hrōzijaną} \textquotedblleft{}to stir\textquotedblright{} (EtymDB).
\end{cousinweb}

\begin{grammarlens}[title={What you've built}]
One more word for telling, explaining and describing at B1.
\end{grammarlens}

\subsection*{Your turn}
\begin{itemize}
\raggedright
  \item \emph{Say it:} \emph{rühren}
  \item \emph{Say it:} \emph{rühren}, once more
  \item \emph{Say it:} a sentence of your own with \emph{rühren}
  \item \emph{Recall:} say the German for to grate, then the German for to halve, then say \emph{rühren} again
\end{itemize}

\begin{tcolorbox}[breakable,colback=teal!4,colframe=teal!35!black,title={Before you move on}]
What does rühren mean? (\textquotedblleft{}To stir; to move (the heart)\textquotedblright{}.) Say it once more.
\end{tcolorbox}
\section[umrühren]{umrühren — to give a stir, to stir round}

\label{lesson:GE-C324-umruehren}



\begin{quote}
Before the new one: say the German for to halve, then the German for to stir.
\end{quote}

\subsection*{You'll want to know: umrühren}
\textbf{umrühren} — \textquotedblleft{}to give a stir, to stir round\textquotedblright{}.

\begin{quote}
Sie rührte die Soße alle fünf Minuten um.
\end{quote}

\begin{quote}
\emph{She stirred the sauce every five minutes.}
\end{quote}

Separable: \emph{er rührt um}.

\begin{cousinweb}
\emph{umrühren} joins \emph{um} \textquotedblleft{}around\textquotedblright{} and \emph{rühren} \textquotedblleft{}to stir\textquotedblright{}, from Proto-Germanic \emph{hrōzijaną} (EtymDB).
\end{cousinweb}

\begin{grammarlens}[title={What you've built}]
One more word for telling, explaining and describing at B1.
\end{grammarlens}

\subsection*{Your turn}
\begin{itemize}
\raggedright
  \item \emph{Say it:} \emph{umrühren}
  \item \emph{Say it:} \emph{umrühren}, once more
  \item \emph{Say it:} a sentence of your own with \emph{umrühren}
  \item \emph{Recall:} say the German for to halve, then the German for to stir, then say \emph{umrühren} again
\end{itemize}

\begin{tcolorbox}[breakable,colback=teal!4,colframe=teal!35!black,title={Before you move on}]
What does umrühren mean? (\textquotedblleft{}To give a stir, to stir round\textquotedblright{}.) Say it once more.
\end{tcolorbox}
\section[verquirlen]{verquirlen — to whisk, to beat (eggs)}

\label{lesson:GE-C324-verquirlen}



\begin{quote}
Before the new one: say the German for to stir, then the German for to give a stir.
\end{quote}

\subsection*{You'll want to know: verquirlen}
\textbf{verquirlen} — \textquotedblleft{}to whisk, to beat (eggs)\textquotedblright{}.

\begin{quote}
Er verquirlte die Eier mit etwas Milch.
\end{quote}

\begin{quote}
\emph{He whisked the eggs with a little milk.}
\end{quote}

Built on \emph{der Quirl}, a kitchen whisk.

\begin{cousinweb}
\emph{verquirlen} is built on \emph{Quirl} \textquotedblleft{}a stirring stick, a whisk\textquotedblright{}, usually traced to Old High German \emph{dweran} \textquotedblleft{}to turn\textquotedblright{}; EtymDB lists no parent for \emph{Quirl}.
\end{cousinweb}

\begin{grammarlens}[title={What you've built}]
One more word for telling, explaining and describing at B1.
\end{grammarlens}

\subsection*{Your turn}
\begin{itemize}
\raggedright
  \item \emph{Say it:} \emph{verquirlen}
  \item \emph{Say it:} \emph{verquirlen}, once more
  \item \emph{Say it:} a sentence of your own with \emph{verquirlen}
  \item \emph{Recall:} say the German for to stir, then the German for to give a stir, then say \emph{verquirlen} again
\end{itemize}

\begin{tcolorbox}[breakable,colback=teal!4,colframe=teal!35!black,title={Before you move on}]
What does verquirlen mean? (\textquotedblleft{}To whisk, to beat (eggs)\textquotedblright{}.) Say it once more.
\end{tcolorbox}
\section[aufschlagen]{aufschlagen — to crack open (an egg); to open (a book)}

\label{lesson:GE-C324-aufschlagen}



\begin{quote}
Before the new one: say the German for to give a stir, then the German for to whisk.
\end{quote}

\subsection*{You'll want to know: aufschlagen}
\textbf{aufschlagen} — \textquotedblleft{}to crack open (an egg); to open (a book)\textquotedblright{}.

\begin{quote}
Sie schlug drei Eier auf.
\end{quote}

\begin{quote}
\emph{She cracked three eggs.}
\end{quote}

Separable: \emph{sie schlägt auf}. Also to open a book: \emph{das Buch aufschlagen}.

\begin{cousinweb}
\emph{aufschlagen} joins \emph{auf} \textquotedblleft{}open\textquotedblright{} and \emph{schlagen} \textquotedblleft{}to strike\textquotedblright{}, an inherited Germanic verb; EtymDB lists no parent for the compound.
\end{cousinweb}

\begin{grammarlens}[title={What you've built}]
One more word for telling, explaining and describing at B1.
\end{grammarlens}

\subsection*{Your turn}
\begin{itemize}
\raggedright
  \item \emph{Say it:} \emph{aufschlagen}
  \item \emph{Say it:} \emph{aufschlagen}, once more
  \item \emph{Say it:} a sentence of your own with \emph{aufschlagen}
  \item \emph{Recall:} say the German for to give a stir, then the German for to whisk, then say \emph{aufschlagen} again
\end{itemize}

\begin{tcolorbox}[breakable,colback=teal!4,colframe=teal!35!black,title={Before you move on}]
What does aufschlagen mean? (\textquotedblleft{}To crack open (an egg); to open (a book)\textquotedblright{}.) Say it once more.
\end{tcolorbox}
\section[der Schneebesen]{der Schneebesen — the whisk}

\label{lesson:GE-C324-schneebesen}



\begin{quote}
Before the new one: say the German for to whisk, then the German for to crack open.
\end{quote}

\subsection*{You'll want to know: der Schneebesen}
\textbf{der Schneebesen} — \textquotedblleft{}the whisk\textquotedblright{}.

\begin{quote}
Mit dem Schneebesen schlug er die Sahne steif.
\end{quote}

\begin{quote}
\emph{He whipped the cream stiff with the whisk.}
\end{quote}

Plural: \emph{die Schneebesen}. \emph{Schnee} 'snow' plus \emph{Besen} 'broom': beaten egg white looks like snow.

\begin{cousinweb}
\emph{Schneebesen} joins \emph{Schnee} \textquotedblleft{}snow\textquotedblright{} and \emph{Besen} \textquotedblleft{}broom\textquotedblright{}, Old High German \emph{besamo}, Proto-Germanic \emph{besamô} (EtymDB).
\end{cousinweb}

\begin{grammarlens}[title={What you've built}]
One more word for telling, explaining and describing at B1.
\end{grammarlens}

\subsection*{Your turn}
\begin{itemize}
\raggedright
  \item \emph{Say it:} \emph{der Schneebesen}
  \item \emph{Say it:} \emph{der Schneebesen}, once more
  \item \emph{Say it:} a sentence of your own with \emph{der Schneebesen}
  \item \emph{Recall:} say the German for to whisk, then the German for to crack open, then say \emph{der Schneebesen} again
\end{itemize}

\begin{tcolorbox}[breakable,colback=teal!4,colframe=teal!35!black,title={Before you move on}]
What does der Schneebesen mean? (\textquotedblleft{}The whisk\textquotedblright{}.) Say it once more.
\end{tcolorbox}
